%% kapa: exam=1 round=24 year=2013 period=2 form=A subject=appraisal_law subject_ko=부동산관계법규 seq=10 number=50
%% source: Q-Net 감정평가사 기출문제 (cst00309) artlSeq=5073802 file=extracted/1008/5073802/2013년 제24회 감정평가사 1차시험 기출문제/2교시(A형).hwp
%% conversion: kordoc 4.24.51 -> tools/gen1.py (auto); review=proofread-sonnet; flags=none
\begin{question}{50}
국토의 계획 및 이용에 관한 법령상 지구단위계획의 내용으로 반드시 포함되어야 하는 것은? (단, 여기서의 지구단위계획은 기존의 용도지구를 폐지하고 그 용도지구에서의 행위제한을 대체하는 사항을 내용으로 하는 지구단위계획이 아니다.)
\begin{choices}
  \item 용도지역이나 용도지구를 대통령령으로 정하는 범위에서 세분하거나 변경하는 사항
  \item 건축물의 배치․형태․색채 또는 건축선에 관한 계획
  \item 계획적인 개발ㆍ정비를 위하여 구획된 일단의 토지의 규모와 조성계획
  \item 건축물의 용도제한, 건축물의 건폐율 또는 용적률, 건축물 높이의 최고한도 또는 최저한도
  \item 간판의 크기․형태․색채 또는 재질
\end{choices}
\end{question}
\answer{4}
